\label{sec:conclusion}

M.8 is the concluding paper in the M-track community character series.
It defines the Composite Community Character Index: a weighted average of
seven community similarity signals derived from federal government data sources,
designed to operationalize the multi-dimensional community bond that
redistricting commissions and courts have sought to preserve since the
communities-of-interest doctrine was first articulated.

The composite does something that no individual M-track signal can do alone:
it provides simultaneous evidence of community cohesion across six distinct
dimensions.
A community that scores well on the composite is not merely economically
coherent or physically coherent or administratively coherent; it is coherent
on \emph{all measured dimensions}, which is the strongest possible argument
for preserving it in a redistricting proceeding.

\paragraph{The court usage template as the primary contribution.}
The court usage template in Section~4 is the primary practical contribution
of this paper.
Each of the seven component signals provides a standalone expert witness
template in its own paper.
M.8 provides the synthesis: the complete expert statement that covers all
seven dimensions at once, the non-partisan declaration that covers all
seven data sources at once, and the brief citation template for practitioners
who need a single sentence reference.

The template is designed to be specific enough to be credible (it cites exact
data sources, exact formulas, exact weight values) and flexible enough to be
useful across states and proceedings (the community identification section
uses fill-in-the-blank language that the expert adapts to the specific
state and specific communities).

\paragraph{Graceful degradation as a design principle.}
The composite formula's most important technical property is graceful
degradation over missing signals.
A practitioner in a rural state without GTFS data and with minimal topographic
variation can still deploy the composite using M.6 + M.1 + M.2 + M.3.
A practitioner in a dense transit city can increase M.7's weight.
A practitioner in a coastal state with significant water boundaries can rely
on M.2's hard-cut rule as an absolute constraint.

The formula normalizes over whatever signals are available, so the resulting
weights always sum to 1.0 over the available components.
This means the composite is always a valid weighted average---there is no
``partial composite'' that needs special handling---and the court statement
does not need to explain why some signals were omitted.

\paragraph{Non-partisanship as an architecture requirement.}
The M.8 composite was designed from the first component (M.0 framework) with
partisan neutrality as a hard architectural requirement.
No component signal uses electoral data.
No weight was chosen to favor a particular partisan outcome.
The calibration argument in Section~3 derives weights entirely from legal
recognition, data quality, and empirical effect considerations.

The partisan neutrality check in Section~5 (Pearson $r < 0.05$ between
composite weights and tract-level Democratic vote fraction) is not an
afterthought; it is the empirical verification of the architectural claim.
If any component signal inadvertently encoded a partisan signal, the
check would detect it.
Passing this check is as important as demonstrating compactness improvement
or within-district variance reduction.

\paragraph{Phase 2 commitments.}
Phase~2 will complete the following quantitative results, all currently
marked ${}^{\dagger}$:
\begin{enumerate}
  \item Within-district economic variance reduction on NC-14.
    Prediction: $\geq$25\% reduction vs.\ geographic baseline.
  \item Zone boundary crossings on NC-14.
    Prediction: zero school district boundary crossings.
  \item Polsby-Popper compactness on NC-14, WI-8, TX-38.
    Prediction: composite $\geq$ geographic.
  \item Partisan neutrality check on NC-14.
    Prediction: $|\rho| < 0.05$.
  \item Partial-data composite validation: NC-14 using M.6 + M.1 + M.2 only
    (three-signal minimum set), verifying graceful degradation.
\end{enumerate}

Phase~2 results will be incorporated into a revised version of this paper
and used to support the M-track series as evidence in redistricting proceedings.
The Phase~2 pipeline run will use:
\begin{verbatim}
bisect state --state NC --year 2020 \
    --weights-override composite-community \
    --version v_m8_phase2
\end{verbatim}

\paragraph{The broader M-track contribution.}
The M-track series (M.0--M.8) is the first systematic operationalization of
communities of interest in algorithmic redistricting.
Prior work in the field has addressed compactness, partisan fairness, racial
equity, and population balance.
Communities of interest---the community type most often cited by residents
in redistricting testimony, and the criterion most directly tied to the
lived experience of people in the affected districts---has been treated as
a qualitative goal that cannot be made quantitative.

The M-track series demonstrates that it can.
Seven distinct community dimensions can be measured, weighted, and combined
into a single composite score for every adjacent tract pair in the country.
That score can be used as an input to a graph partitioning algorithm that
respects communities automatically, without requiring anyone to draw lines
by hand.
And the composite score can be presented to courts and commissions as
specific, quantitative, non-partisan evidence that the submitted plan
preserves the communities it claims to preserve.

This is the promise of algorithmic redistricting: not that algorithms make
better decisions than humans, but that they make decisions that can be
explained, verified, and defended.
The composite community character index is that explanation.
